\documentclass[a4paper,11pt]{article}
\usepackage[utf8]{inputenc}
\usepackage[T1]{fontenc}
\usepackage[french]{babel}
\usepackage{lmodern}
\usepackage{amsmath,amssymb,amsthm}
\usepackage{mathrsfs}
\usepackage{enumitem}
\usepackage{multicol}

\setlist[enumerate]{label=\textbf{\textcolor{Couleur8}{\arabic*.}}, left=1em}

% Si vous voulez aussi modifier les listes enumerate imbriquées :
\setlist[enumerate,2]{label=\textcolor{Couleur8}{\alph*)}}
\setlist[enumerate,3]{label=\textcolor{Couleur8}{\roman*)}}
\setlist[enumerate,4]{label=\textcolor{Couleur8}{\Alph*)}}
\usepackage{graphicx}
\usepackage{xcolor}
\usepackage{tcolorbox}
\usepackage{geometry}
\usepackage{fancyhdr}
\usepackage{titlesec}
\usepackage{cancel}
\usepackage{ifthen}
\usepackage{tikz}
\usetikzlibrary{arrows.meta}
\usepackage{tkz-tab}
\usepackage{eurosym}

% OPTION PROF/ÉLÈVE
% Décommentez UNE des deux lignes suivantes :
\newboolean{prof}\setboolean{prof}{true}   % Version professeur (avec corrections des devoirs)
%\newboolean{prof}\setboolean{prof}{false}  % Version élève (sans corrections des devoirs)

\definecolor{Couleur1}{HTML}{4A5D7A} %Th
\definecolor{Couleur2}{HTML}{A5B5D6} %Prop
\definecolor{Couleur3}{HTML}{A8C68A} %Méthode
\definecolor{Couleur6}{HTML}{8A9CC1} %Def
\definecolor{Couleur7}{HTML}{E6B459} %Rq Voc
\definecolor{Couleur8}{HTML}{D36740} %Titres
\definecolor{correction}{HTML}{D36740} % Couleur pour les corrections
\definecolor{coopmathscorr}{HTML}{F15929} % Couleur corrections coopmaths

% Configuration des marges
\geometry{
	top=2cm,
	bottom=2cm,
	inner=1.5cm,
	outer=1.5cm,
}

% Police
\renewcommand{\familydefault}{\sfdefault}

% Compteurs
\newcounter{seancenum}
\newcounter{seancenumD}
\setcounter{seancenum}{1}
\setcounter{seancenumD}{1}

% Configuration des en-têtes et pieds de page
\pagestyle{fancy}
\fancyhf{}
\ifthenelse{\boolean{prof}}{
    \fancyhead[L]{\textcolor{Couleur8}{\textbf{Corrections Automatismes - Classe de seconde - Version Professeur}}}
}{
    \fancyhead[L]{\textcolor{Couleur8}{\textbf{Corrections Automatismes - Séance 1 - Classe de seconde}}}
}
\fancyhead[R]{\textcolor{Couleur8}{\textbf{2025-2026}}}
\fancyfoot[L]{\textcolor{Couleur8}{Mme Bertail et Mme Pinto}}
\fancyfoot[R]{\textcolor{Couleur8}{\thepage}}
\renewcommand{\headrulewidth}{1pt}
\renewcommand{\headrule}{\hbox to\headwidth{\color{Couleur8}\leaders\hrule height \headrulewidth\hfill}}

% Environnements personnalisés
\newtcolorbox{seance}[1]{
    colback=Couleur6!10,
    colframe=Couleur1!65,
    fonttitle=\bfseries\large,
    title=Correction Séance \theseancenum #1,
    left=5mm,
    right=5mm,
    top=3mm,
    bottom=3mm,
    arc=0mm,
    after=\stepcounter{seancenum}
}

\newtcolorbox{seanceD}[1]{
    colback=Couleur6!5,
    colframe=Couleur6!45,
    fonttitle=\bfseries\large,
    title=Correction Devoirs - Séance \theseancenumD #1,
    left=5mm,
    right=5mm,
    top=3mm,
    bottom=3mm,
    arc=0mm,
    after=\stepcounter{seancenumD}
}

\begin{document}

\setcounter{seancenum}{1}
\newpage
\begin{seance}{} %1

\begin{enumerate}
\item $\dfrac{2}{6}-\dfrac{3}{9}=\dfrac{2{\color{Couleur1}\boldsymbol{\times 3}}}{6{\color{Couleur1}\boldsymbol{\times 3}}}-\dfrac{3{\color{Couleur1}\boldsymbol{\times 2}}}{9{\color{Couleur1}\boldsymbol{\times 2}}}=\dfrac{6-6}{18}=\dfrac{0}{18}=$ ${\color{correction}\boldsymbol{0}}$

On a réduit les fractions au même dénominateur (18), puis effectué la soustraction.

\item $\dfrac{8}{6}-\dfrac{2}{21}=\dfrac{8{\color{Couleur1}\boldsymbol{\times 7}}}{6{\color{Couleur1}\boldsymbol{\times 7}}}-\dfrac{2{\color{Couleur1}\boldsymbol{\times 2}}}{21{\color{Couleur1}\boldsymbol{\times 2}}}=\dfrac{56-4}{42}=\dfrac{52}{42}=\dfrac{26\times2}{21\times2}=$ ${\color{correction}\boldsymbol{\dfrac{26}{21}}}$

Le dénominateur commun est 42. On simplifie ensuite par 2.

\item $\dfrac{11}{10} \div \dfrac{77}{-30} = \dfrac{11}{10} \times \dfrac{-30}{77} = \dfrac{11 \times (-30)}{10 \times 77} = \dfrac{-330}{770}$

On décompose : $\dfrac{(-11) \times 30}{10 \times 77} = \dfrac{(-11) \times 3 \times 10}{10 \times 7 \times 11} = \dfrac{-3 \times \cancel{10} \times \cancel{11}}{\cancel{10} \times 7 \times \cancel{11}} = {\color{correction}\boldsymbol{\dfrac{3}{7}}}$

Pour diviser par une fraction, on multiplie par son inverse.

\item $\dfrac{-15}{55}\times\dfrac{-77}{40}=\dfrac{(-15)\times(-77)}{55\times40} = \dfrac{15 \times 77}{55 \times 40}$

On décompose : $\dfrac{3 \times 5 \times 7 \times 11}{5 \times 11 \times 8 \times 5} = \dfrac{3 \times \cancel{5} \times 7 \times \cancel{11}}{8 \times 5 \times \cancel{5} \times \cancel{11}} = {\color{correction}\boldsymbol{\dfrac{21}{40}}}$

Le produit de deux nombres négatifs est positif.

\item Le matin, Aude a bu $\dfrac{5}{7}$ de la bouteille. Il reste alors $1 - \dfrac{5}{7} = \dfrac{7}{7} - \dfrac{5}{7} = \dfrac{2}{7}$ de la bouteille.

À midi, elle a bu $\dfrac{1}{6}$ du reste, soit : $\dfrac{1}{6} \times \dfrac{2}{7} = \dfrac{2}{42} = {\color{correction}\boldsymbol{\dfrac{1}{21}}}$ de la bouteille.

\item $A = \dfrac{1}{2-\dfrac{1}{2}} = \dfrac{1}{\dfrac{4}{2}-\dfrac{1}{2}} = \dfrac{1}{\dfrac{3}{2}} = 1 \times \dfrac{2}{3} = {\color{correction}\boldsymbol{\dfrac{2}{3}}}$

\end{enumerate}

\end{seance}

\end{document}